% sec_gram_space.tex -- Mode Gram kernel, RKHS completion, and centred reconstruction
\section{Mode Gram Space and Centred Reconstruction}%
\label{sec:gram-space}

This appendix builds the Hilbert coordinate system used by the Poisson strip
observation projection.  The secant Gram kernel, the unitary RKHS completion,
and the centred Poisson reconstruction statement identify the observation
channel later projected onto the integer-orbit strip subspace.  Standard RKHS
completion and Laplace/characteristic-function uniqueness are used as cited
background.

\subsection{The mode Gram kernel and Hilbert completion}

Let
\[
\omega(\dd y):=\frac{\dd y}{\pi(1+y^2)}.
\]
For $\varepsilon\in\RR$ define the secant family
\begin{equation}\label{eq:mode-secant-family}
\Psi_\varepsilon(y):=\frac{R_\varepsilon(y)-1}{\varepsilon}
=\frac{2y-\varepsilon}{1+(y-\varepsilon)^2},
\qquad y\in\RR,
\end{equation}
where the formula remains valid at $\varepsilon=0$ and gives
\[
\Psi_0(y)=u_1(y)=\frac{2y}{1+y^2}.
\]
Since
\[
R_\varepsilon(y)\,\omega(\dd y)=P_1(y-\varepsilon)\,\dd y,
\]
one has
\[
\int_{\RR}R_\varepsilon(y)\,\omega(\dd y)=1
\qquad(\varepsilon\in\RR),
\]
and therefore $\Psi_\varepsilon\in L^2_0(\omega)$ for every
$\varepsilon\in\RR$.

\begin{theorem}[The mode Gram kernel]%
\label{thm:mode-gram-kernel}
For every $\varepsilon,\eta\in\RR$,
\begin{equation}\label{eq:mode-gram-kernel}
\langle \Psi_\varepsilon,\Psi_\eta\rangle_{L^2(\omega)}
=\frac{2}{4+(\varepsilon-\eta)^2}.
\end{equation}
Equivalently,
\begin{equation}\label{eq:mode-gram-RR}
\int_{\RR}R_\varepsilon(y)R_\eta(y)\,\omega(\dd y)
=\frac{\varepsilon^2+\eta^2+4}{(\varepsilon-\eta)^2+4}.
\end{equation}
\end{theorem}

\begin{proof}
For $\varepsilon\neq\eta$ define
\[
I(\varepsilon,\eta)
:=\frac{1}{\pi}\int_{\RR}
\frac{1+y^2}{(1+(y-\varepsilon)^2)(1+(y-\eta)^2)}\,\dd y.
\]
The integrand is $O(|z|^{-2})$ on the upper half-plane and has simple
poles at $z=\varepsilon+\ii$ and $z=\eta+\ii$.  A direct residue
computation gives
\[
\operatorname{Res}_{z=\varepsilon+\ii}
=\frac{\varepsilon(\varepsilon+2\ii)}
{2\ii(\varepsilon-\eta)(\varepsilon-\eta+2\ii)},
\]
\[
\operatorname{Res}_{z=\eta+\ii}
=\frac{\eta(\eta+2\ii)}
{2\ii(\eta-\varepsilon)(\eta-\varepsilon+2\ii)}.
\]
Hence, by the residue theorem,
\[
I(\varepsilon,\eta)
=2\ii\Bigl(
\operatorname{Res}_{z=\varepsilon+\ii}
+\operatorname{Res}_{z=\eta+\ii}
\Bigr)
=\frac{\varepsilon^2+\eta^2+4}{(\varepsilon-\eta)^2+4}.
\]
This proves \eqref{eq:mode-gram-RR} for $\varepsilon\neq\eta$, and the
diagonal case follows by continuity.

Since
\[
\int_{\RR}R_\varepsilon(y)\,\omega(\dd y)=1,
\]
one obtains
\[
\int_{\RR}(R_\varepsilon(y)-1)(R_\eta(y)-1)\,\omega(\dd y)
=\frac{2\varepsilon\eta}{4+(\varepsilon-\eta)^2}.
\]
For $\varepsilon\eta\neq 0$, division by $\varepsilon\eta$ yields
\eqref{eq:mode-gram-kernel}; the remaining cases follow by continuity in
$(\varepsilon,\eta)$.
\end{proof}

\begin{corollary}[Closed formulas for quadratic mode integrals]%
\label{cor:mode-quadratic-integrals}
For $|\varepsilon|<1$ one has the $L^2(\omega)$ expansion
\[
\Psi_\varepsilon
=\sum_{n=1}^{\infty}u_n\,\varepsilon^{n-1}.
\]
Consequently, for every $p,q\ge 1$,
\[
\langle u_{2p},u_{2q}\rangle_{L^2(\omega)}
=(-1)^{p+q}2^{-(2p+2q-1)}
\binom{2p+2q-2}{2p-1},
\]
for every $p,q\ge 0$,
\[
\langle u_{2p+1},u_{2q+1}\rangle_{L^2(\omega)}
=(-1)^{p+q}2^{-(2p+2q+1)}
\binom{2p+2q}{2p},
\]
and
\[
\langle u_{2p},u_{2q+1}\rangle_{L^2(\omega)}=0.
\]
\end{corollary}

\begin{proof}
Use $y=\tan\theta$ to identify $L^2(\omega)$ with
$L^2((-\pi/2,\pi/2),\dd\theta/\pi)$.  From
Theorem~\ref{thm:cayley-chebyshev-mode},
\[
u_n(\tan\theta)=\cos^n\theta\,U_n(\sin\theta).
\]
For $s\in[-1,1]$ the Chebyshev polynomials of the second kind satisfy
$|U_n(s)|\le n+1$.  Hence, for every $0<r<1$ and
$|\varepsilon|\le r$,
\[
\sum_{n=1}^{\infty}
\sup_{\theta}\bigl|u_n(\tan\theta)\varepsilon^{n-1}\bigr|
\le
\sum_{n=1}^{\infty}(n+1)r^{n-1}<\infty.
\]
The Weierstrass test gives uniform, and therefore $L^2(\omega)$,
convergence.  The Chebyshev generating function used in
Theorem~\ref{thm:cayley-chebyshev-mode} is valid because
$|\varepsilon\cos\theta|<1$, so the uniform limit is
$(R_\varepsilon-1)/\varepsilon=\Psi_\varepsilon$.

Moreover $\|u_n\|_{L^2(\omega)}\le n+1$, and therefore, for
$|\varepsilon|,|\eta|\le r<1$,
\[
\sum_{m,n\ge1}
\|u_m\|_{L^2(\omega)}\|u_n\|_{L^2(\omega)}
|\varepsilon|^{m-1}|\eta|^{n-1}<\infty.
\]
Thus the inner product of the two $L^2$ series may be expanded
absolutely:
\[
\frac{2}{4+(\varepsilon-\eta)^2}
=\sum_{m,n\ge 1}
\langle u_m,u_n\rangle_{L^2(\omega)}\,
\varepsilon^{m-1}\eta^{n-1}.
\]
On the other hand,
\[
\frac{2}{4+(\varepsilon-\eta)^2}
=\frac12\sum_{\ell=0}^{\infty}
(-1)^\ell 4^{-\ell}(\varepsilon-\eta)^{2\ell}.
\]
Comparing coefficients gives the stated formulas.  Mixed parity terms
vanish because the kernel only contains even total degree.
\end{proof}

\begin{lemma}[Poisson injectivity in the Cauchy-weighted space]%
\label{lem:poisson-injectivity-weighted}
Let
\[
\omega(\dd y)=\frac{\dd y}{\pi(1+y^2)}
\]
and let $f\in L^2(\omega)$.  If
\[
P_1*f=0
\qquad\text{in }\mathscr S'(\RR),
\]
then $f=0$ in $L^2(\omega)$.  Equivalently, convolution by $P_1$ is
injective on the Cauchy-weighted space.
\end{lemma}

\begin{proof}
Since $f\in L^2(\omega)$,
\[
\left|\int_{\RR} f(y)\varphi(y)\dd y\right|
\le
\left(\int_{\RR}\frac{|f(y)|^2}{1+y^2}\dd y\right)^{1/2}
\left(\int_{\RR}(1+y^2)|\varphi(y)|^2\dd y\right)^{1/2}
\]
for every $\varphi\in\mathscr S(\RR)$.  Thus $f$ defines a tempered
distribution.  Fourier transforming the distributional identity \(P_1*f=0\)
gives
\[
e^{-|\xi|}\widehat f(\xi)=0
\qquad\text{in }\mathscr S'(\RR).
\]
On \(\RR\setminus\{0\}\) the multiplier \(e^{-|\xi|}\) is smooth and
strictly positive.  Hence, for every
\(\phi\in C_c^\infty(\RR\setminus\{0\})\), the test function
\(\psi=e^{|\xi|}\phi\) is again smooth and compactly supported away from the
origin, and
\[
\langle\widehat f,\phi\rangle
=\langle e^{-|\xi|}\widehat f,\psi\rangle=0 .
\]
Therefore \(\operatorname{supp}\widehat f\subset\{0\}\).  By the standard
structure theorem for distributions supported at one point
\cite{Grafakos2014}, there are constants \(a_0,\ldots,a_M\) such that
\[
\widehat f=\sum_{j=0}^{M}a_j\delta_0^{(j)}.
\]
After inverse Fourier transform, \(f\) is a polynomial distribution.  Since
the same \(f\) is represented by an element of \(L^2(\omega)\), that
polynomial has a locally integrable polynomial representative.  The
Cauchy-weighted condition excludes every nonconstant polynomial: if
\(p(y)\) has degree \(d\ge1\), then
\(|p(y)|^2(1+y^2)^{-1}\) is comparable to \(|y|^{2d-2}\) along large
\(|y|\), which is not integrable.  Hence \(f=c\) a.e. for some constant
\(c\).  Finally \(P_1*c=c\) in \(\mathscr S'(\RR)\), because \(P_1\) has
total mass one.  The hypothesis \(P_1*f=0\) forces \(c=0\), so \(f=0\)
a.e.
\end{proof}

\begin{theorem}[The mode space and its RKHS completion]%
\label{thm:mode-space-rkhs}
Let
\[
L^2_0(\omega):=
\left\{
f\in L^2(\omega):\ \int_{\RR}f(y)\,\omega(\dd y)=0
\right\},
\]
and let $\mathcal H_K$ be the reproducing-kernel Hilbert space generated
by the kernel
\[
K(a,b):=\frac{2}{4+(a-b)^2}.
\]
Set
\[
\mathcal S_0:=\operatorname{span}\{\Psi_\varepsilon:\varepsilon\in\RR\},
\qquad
\mathcal S:=\overline{\mathcal S_0}^{\,L^2(\omega)}
\subset L^2_0(\omega),
\]
Then
\[
\mathcal S=L^2_0(\omega).
\]
Equivalently, the algebraic span
$\mathcal S_0$ is dense in~$L^2_0(\omega)$.
The density assertion rests on the isolated Poisson injectivity lemma
immediately above: \(P_1*f=0\) first forces
\(\operatorname{supp}\widehat f\subset\{0\}\), and the Cauchy-weighted
\(L^2(\omega)\) condition excludes the resulting polynomial distributions.
Moreover, the linear assignment
\[
\mathcal U_0:\mathcal S_0\longrightarrow \mathcal H_K,
\qquad
\mathcal U_0(\Psi_\varepsilon):=K(\cdot,\varepsilon)
\qquad(\varepsilon\in\RR)
\]
is an isometry and extends uniquely to a unitary operator
\[
\mathcal U:L^2_0(\omega)\longrightarrow \mathcal H_K.
\]
\end{theorem}

\begin{proof}
We already noted that $\Psi_\varepsilon\in L^2_0(\omega)$, so
$\mathcal S_0\subset L^2_0(\omega)$.  Let
$f\in L^2_0(\omega)$ be orthogonal to~$\mathcal S_0$.  Then for every
$\varepsilon\in\RR$,
\[
0=\langle f,\Psi_\varepsilon\rangle_{L^2(\omega)}.
\]
If $\varepsilon\neq 0$, multiplying by~$\varepsilon$ gives
\[
0=\int_{\RR}f(y)\bigl(R_\varepsilon(y)-1\bigr)\,\omega(\dd y)
=\int_{\RR}f(y)R_\varepsilon(y)\,\omega(\dd y),
\]
because $\int f\,\dd\omega=0$.  Since
$R_\varepsilon(y)\,\omega(\dd y)=P_1(y-\varepsilon)\,\dd y$, this means
\[
0=\int_{\RR}f(y)P_1(\varepsilon-y)\,\dd y
\qquad(\varepsilon\in\RR).
\]
At $\varepsilon=0$ this is exactly the zero-mean condition.
The weighted Cauchy--Schwarz estimate used in
Lemma~\ref{lem:poisson-injectivity-weighted} shows that these pointwise
integrals are finite.  It also identifies the pointwise convolution with the
distributional one: for $\varphi\in\mathscr S(\RR)$,
\[
\int_{\RR}\left(\int_{\RR}f(y)P_1(\varepsilon-y)\dd y\right)
\varphi(\varepsilon)\dd\varepsilon
=\langle f,\widetilde P_1*\varphi\rangle,
\qquad \widetilde P_1(x):=P_1(-x)=P_1(x),
\]
with Fubini justified because $\widetilde P_1*\varphi(y)=O((1+|y|)^{-2})$
at infinity.  Thus the distributional convolution $P_1*f$ vanishes.
Lemma~\ref{lem:poisson-injectivity-weighted} applies and gives $f=0$.
Therefore $\mathcal S_0^\perp=\{0\}$ in $L^2_0(\omega)$, so
\[
\overline{\mathcal S_0}^{\,L^2(\omega)}=L^2_0(\omega).
\]
By definition this closure is~$\mathcal S$, hence
$\mathcal S=L^2_0(\omega)$.

Theorem~\ref{thm:mode-gram-kernel} shows that
\[
\langle \Psi_\varepsilon,\Psi_\eta\rangle_{L^2(\omega)}
=K(\varepsilon,\eta)
=\langle K(\cdot,\varepsilon),K(\cdot,\eta)\rangle_{\mathcal H_K},
\]
so $\mathcal U_0$ is an isometry on~$\mathcal S_0$.  Since the
kernel sections span a dense subspace of~$\mathcal H_K$, the isometry
extends uniquely to a surjective isometry
$\mathcal U:L^2_0(\omega)\to\mathcal H_K$, that is, to a unitary
operator.
\end{proof}

\subsection{Centred reconstruction from Poisson data}

\begin{proposition}[Centred Poisson reconstruction by Laplace uniqueness]%
\label{prop:poisson-central-two-channel}
Let $\nu$ be a probability measure on~$\RR$ with finite mean~$\bar\gamma$.
Let $\nu_c$ denote the law of the centred variable
$\Delta:=\gamma-\bar\gamma$.  Define
\[
h_t(x):=\int_{\RR}P_t(x-\gamma)\dd\nu(\gamma),
\]
the central observation functions
\[
A(t):=\pi\,h_t(\bar\gamma),
\qquad
B(t):=\pi\,\partial_x h_t(x)\big|_{x=\bar\gamma}
\qquad(t>0).
\]
\[
H(t):=\int_t^\infty B(s)\dd s
\qquad(t>0),
\]
and let
\[
\phi_{\nu_c}(r):=\int_{\RR}e^{ir\Delta}\dd\nu(\gamma),
\qquad r\in\RR,
\]
be the centred characteristic function.  Then for every $t>0$ one has
\begin{equation}\label{eq:poisson-laplace-reconstruction}
A(t)+\ii H(t)
=\int_0^\infty e^{-tr}\phi_{\nu_c}(r)\dd r.
\end{equation}
Consequently, the pair $(A,B)$ determines the centred law~$\nu_c$
uniquely.  Equivalently, the triple $(\bar\gamma,A,B)$ determines the
original measure~$\nu$ uniquely.
\end{proposition}

\begin{proof}
By the elementary Laplace identities
\[
\int_0^\infty e^{-tr}\cos(r\Delta)\dd r=\frac{t}{t^2+\Delta^2},
\qquad
\int_0^\infty e^{-tr}\sin(r\Delta)\dd r=\frac{\Delta}{t^2+\Delta^2},
\]
one obtains
\[
A(t)
=\int_{\RR}\frac{t}{t^2+\Delta^2}\dd\nu(\gamma)
=\int_0^\infty e^{-tr}\Re\phi_{\nu_c}(r)\dd r.
\]
Differentiation under the integral is justified by dominated convergence,
since for each fixed $t>0$,
\[
\left|\frac{2t(x-\gamma)}{((x-\gamma)^2+t^2)^2}\right|
\le \frac{9}{8\sqrt3}\,t^{-2}
\qquad(x,\gamma\in\RR),
\]
and
\[
\partial_x h_t(x)\big|_{x=\bar\gamma}
=\frac{2}{\pi}\int_{\RR}\frac{t\,\Delta}{(\Delta^2+t^2)^2}\dd\nu(\gamma),
\]
in particular
\[
B(t)=2t\int_{\RR}\frac{\Delta}{(\Delta^2+t^2)^2}\dd\nu(\gamma).
\]
Since
\[
\int_t^\infty \frac{2s|\Delta|}{(\Delta^2+s^2)^2}\dd s
=\frac{|\Delta|}{\Delta^2+t^2}
\le \frac{1}{2t},
\]
Fubini's theorem applies and gives
\[
H(t)
=\int_{\RR}\frac{\Delta}{t^2+\Delta^2}\dd\nu(\gamma)
=\int_0^\infty e^{-tr}\Im\phi_{\nu_c}(r)\dd r.
\]
Combining the two formulas yields
\eqref{eq:poisson-laplace-reconstruction}.  Let
\[
F(t):=A(t)+\ii H(t)
=\int_0^\infty e^{-tr}\phi_{\nu_c}(r)\dd r,
\qquad t>0 .
\]
By uniqueness of the Laplace transform on bounded continuous functions
\cite{Widder1941}, \(F\) determines \(\phi_{\nu_c}(r)\) for every
\(r\ge0\).  Since \(\nu_c\) is a probability measure on \(\RR\), its
characteristic function satisfies the conjugate symmetry
\[
\phi_{\nu_c}(-r)=\overline{\phi_{\nu_c}(r)}
\qquad(r\ge0).
\]
Thus the values on \( [0,\infty) \) determine the values on all of
\(\RR\).  Levy's continuity/inversion theorem for characteristic functions
then recovers the centred law~$\nu_c$ \cite{Lukacs1970,Feller1971}, and
adding the mean recovers~$\nu$.
\end{proof}

\begin{corollary}[Symmetric channel]%
\label{cor:poisson-symmetrization}
Under the hypotheses of Proposition~\ref{prop:poisson-central-two-channel}, the
function $A$ determines the symmetrization of the centred law~$\nu_c$.
\end{corollary}

\begin{proof}
The real part of $\phi_{\nu_c}$ is exactly the characteristic function of
the centred symmetrization
$\frac12(\nu_c+\check\nu_c)$, where $\check\nu_c(E):=\nu_c(-E)$.  Since
$A$ is the Laplace transform of $\Re\phi_{\nu_c}$ on $[0,\infty)$, Laplace
inversion recovers that real part on the positive half-line.  Conjugate
symmetry makes $\Re\phi_{\nu_c}$ even, so this determines the whole
characteristic function of the symmetrized law.  This is the even channel.
\end{proof}

\begin{corollary}[Symmetric case from the odd channel]%
\label{cor:poisson-symmetric-case}
Under the hypotheses of Proposition~\ref{prop:poisson-central-two-channel}, the
centred law~$\nu_c$ is symmetric if and only if $B\equiv 0$.  In that case
$A$ alone determines~$\nu_c$, and the pair $(\bar\gamma,A)$ determines the
original law~$\nu$.
\end{corollary}

\begin{proof}
The centred law is symmetric if and only if
$\Im\phi_{\nu_c}\equiv 0$ on~$\RR$.  By
\eqref{eq:poisson-laplace-reconstruction} and Laplace uniqueness,
\(B\equiv0\) is equivalent to \(H\equiv0\), hence to
\(\Im\phi_{\nu_c}\equiv0\) on \([0,\infty)\).  Conjugate symmetry then makes
\(\Im\phi_{\nu_c}\equiv0\) on all of \(\RR\), which is equivalent to
symmetry of~\(\nu_c\).  This is the odd channel obstruction: when it
vanishes, the centred law agrees with its symmetrization.  Therefore
Corollary~\ref{cor:poisson-symmetrization} shows that $A$ determines~$\nu_c$.
Translating by the known mean $\bar\gamma$ then determines~$\nu$.
\end{proof}

\begin{theorem}[Observation channels as evaluation functionals]%
\label{thm:poisson-channel-rkhs}
Let $\nu$ be a probability measure on~$\RR$ with finite mean~$\bar\gamma$.
For $t>0$ define the centred fluctuation profile
\[
\mathfrak d_t(y):=
\frac{h_t(\bar\gamma+ty)}{g_t(\bar\gamma+ty)}-1,
\qquad y\in\RR.
\]
Then $\mathfrak d_t\in L^2_0(\omega)$ and
\begin{equation}\label{eq:delta-t-secant}
\mathfrak d_t(y)
=\frac{1}{t}\int_{\RR}\Delta\,\Psi_{\Delta/t}(y)\,\dd\nu(\gamma).
\end{equation}
Consequently,
\begin{equation}\label{eq:delta-t-rkhs}
(\mathcal U\mathfrak d_t)(a)
=\frac{1}{t}\int_{\RR}\Delta\,K(a,\Delta/t)\,\dd\nu(\gamma),
\qquad a\in\RR.
\end{equation}
Moreover,
\begin{equation}\label{eq:delta-t-origin}
\mathfrak d_t(0)=tA(t)-1,
\end{equation}
and
\begin{equation}\label{eq:delta-t-rkhs-origin}
(\mathcal U\mathfrak d_t)(0)=2t\,H(2t)
=2t\int_{2t}^{\infty}B(s)\,\dd s.
\end{equation}
\end{theorem}

\begin{proof}
By definition,
\[
\mathfrak d_t(y)
=\int_{\RR}\bigl(R_{\Delta/t}(y)-1\bigr)\,\dd\nu(\gamma)
=\frac{1}{t}\int_{\RR}\Delta\,\Psi_{\Delta/t}(y)\,\dd\nu(\gamma),
\]
which is \eqref{eq:delta-t-secant}.  By Theorem~\ref{thm:mode-gram-kernel},
\[
\|\Psi_\varepsilon\|_{L^2(\omega)}^2
=\frac{2}{4+(\varepsilon-\varepsilon)^2}
=\frac12.
\]
Hence
\[
\left\|\mathfrak d_t\right\|_{L^2(\omega)}
\le\frac{1}{t}\int_{\RR}|\Delta|\,\left\|\Psi_{\Delta/t}\right\|_{L^2(\omega)}\,\dd\nu(\gamma)
=\frac{1}{\sqrt2\,t}\int_{\RR}|\Delta|\,\dd\nu(\gamma)
<\infty,
\]
so \eqref{eq:delta-t-secant} is a Bochner integral in $L^2_0(\omega)$.
Therefore,
\[
\int_{\RR}\mathfrak d_t(y)\,\omega(\dd y)
=\int_{\RR}(h_t(x)-g_t(x))\,\dd x=0,
\]
and thus $\mathfrak d_t\in L^2_0(\omega)$.  Applying the unitary map
$\mathcal U$ from Theorem~\ref{thm:mode-space-rkhs} to
\eqref{eq:delta-t-secant} yields \eqref{eq:delta-t-rkhs}.

At $y=0$ one has
\[
g_t(\bar\gamma)=\frac{1}{\pi t},
\]
so
\[
\mathfrak d_t(0)
=\frac{h_t(\bar\gamma)}{g_t(\bar\gamma)}-1
=tA(t)-1,
\]
which proves \eqref{eq:delta-t-origin}.  Finally, evaluating
\eqref{eq:delta-t-rkhs} at $a=0$ and using
$K(0,b)=2/(4+b^2)$ gives
\[
(\mathcal U\mathfrak d_t)(0)
=\frac{1}{t}\int_{\RR}\Delta\frac{2}{4+(\Delta/t)^2}\,\dd\nu(\gamma)
=2t\int_{\RR}\frac{\Delta}{4t^2+\Delta^2}\,\dd\nu(\gamma).
\]
The definition of $H$ in
Proposition~\ref{prop:poisson-central-two-channel} gives
\[
H(2t)=\int_{\RR}\frac{\Delta}{4t^2+\Delta^2}\,\dd\nu(\gamma),
\]
and \eqref{eq:delta-t-rkhs-origin} follows.
\end{proof}

\begin{remark}
The previous theorem recovers the centred Cauchy transform on the
imaginary axis:
\[
G_{\nu_c}(\ii t)
=\int_{\RR}\frac{1}{\ii t-\Delta}\dd\nu(\gamma)
=-H(t)-\ii A(t)
\qquad(t>0).
\]
Thus the Laplace-transform formulation strengthens the earlier
Cauchy-transform identity while also isolating the loss of the translation
parameter.  In particular, the proof of
Proposition~\ref{prop:poisson-central-two-channel} is better viewed as a
Laplace-uniqueness plus characteristic-function-inversion argument than
as a Stieltjes inversion argument.
\end{remark}
